\documentclass[a4paper,11pt]{article}
\usepackage[utf8]{inputenc}
\usepackage[T1]{fontenc}
\usepackage[ngerman]{babel}
\usepackage{amsmath,amssymb}
\usepackage{geometry}
\geometry{margin=2.2cm}
\usepackage{enumitem}
\usepackage{fancyhdr}
\usepackage{xcolor}

\definecolor{base}{HTML}{1E1E2E}
\definecolor{fg}{HTML}{CDD6F4}
\definecolor{accent}{HTML}{89B4FA}
\definecolor{rulecol}{HTML}{45475A}
\pagecolor{base}

\newcommand{\transp}{^{\mathsf{T}}}
\newcommand{\norm}[1]{\left\lVert #1 \right\rVert}
\newcommand{\antwort}[1]{\par\vspace{#1}}

\usepackage{titlesec}
\titleformat{\section}{\color{accent}\normalfont\large\bfseries}{}{0pt}{}

\pagestyle{fancy}\fancyhf{}
\lhead{\color{fg}\small Abfrage -- Matrizen}
\rhead{\color{fg}\small Selbsttest 01}
\cfoot{\color{fg}\thepage}
\renewcommand{\headrule}{{\color{rulecol}\hrule height 0.4pt}}
\setlength{\parindent}{0pt}

\begin{document}
\color{fg}
\begin{center}
  {\color{accent}\LARGE\bfseries Selbsttest 01 -- deine Schwachstellen}\\[0.2em]
  {\small Erst selbst überlegen, dann mit \texttt{results/erklaerung/abfrage\_01\_loesung.pdf} kontrollieren.}
\end{center}
\vspace{0.4em}{\color{rulecol}\hrule}\vspace{0.6em}

\section*{Verständnis / Wahr--Falsch}
\begin{enumerate}[label=\arabic*.,leftmargin=*]
  \item Wahr oder falsch: Eine symmetrische positiv definite Matrix kann einen Eigenwert $0$ besitzen. Begründe.
  \item Wahr oder falsch: Jede quadratische Matrix ist diagonalisierbar. Wann genau geht es?
  \item Wann ist eine Hyperebene $\vec n\transp\vec x + n_0=0$ ein Untervektorraum?
  \item Welche Dimensionen haben $U,\Sigma,V$ in der SVD einer $m\times n$-Matrix?
        Was steht in den Spalten von $V$?
  \item Bei der Cholesky-Zerlegung: durch welche Größe teilst du, um $\ell_{ij}$ ($i>j$) zu bestimmen?
\end{enumerate}
\antwort{4.5cm}

\section*{Rechnen}
\begin{enumerate}[label=\arabic*.,leftmargin=*,resume]
  \item Ist $A=\begin{pmatrix}2&1\\1&2\end{pmatrix}$ positiv definit?
        Prüfe über die führenden Hauptminoren \emph{und} über die Eigenwerte.
  \item Bestimme die Inverse von $\begin{pmatrix}4&3\\2&2\end{pmatrix}$ mit der 2$\times$2-Formel.
  \item Cholesky von $A=\begin{pmatrix}9&3\\3&5\end{pmatrix}$: bestimme $L$ mit $A=LL\transp$.
  \item $E:\ x_1+2x_2+2x_3=6$. Abstand des Punktes $p=(1,1,1)\transp$ zu $E$?
  \item SVD-Start: für $A=\begin{pmatrix}3&0\\0&0\\0&4\end{pmatrix}$ berechne $A\transp A$ und die Singulärwerte.
\end{enumerate}
\antwort{6.5cm}

\section*{Erklären}
\begin{enumerate}[label=\arabic*.,leftmargin=*,resume]
  \item Leite für $-u''=f$ mit dem zentralen Differenzenquotienten die Gleichung
        $-u_{i-1}+2u_i-u_{i+1}=h^2 f_i$ her. Warum ist die Systemmatrix tridiagonal?
  \item Warum löst man bei linearer Regression $X\transp X\,\vec c=X\transp\vec y$
        statt $X\vec c=\vec y$?
\end{enumerate}
\antwort{4cm}

\end{document}
